\begin{tikzpicture}[
  >=stealth,
  arrow/.style={-{Stealth[length=2.5mm]}, thick, color=black!80},
  arrow_dash/.style={-{Stealth[length=2.5mm]}, thick, dashed, color=black!60},
  data/.style={rectangle, draw, fill=blue!10, minimum width=1.8cm, minimum height=0.9cm, align=center, font=\small\sffamily, rounded corners=2pt, draw=blue!60!black},
  core/.style={rectangle, draw, rounded corners=4pt, fill=orange!30, minimum width=2.2cm, minimum height=1.2cm, align=center, font=\small\sffamily\bfseries, draw=orange!80!black, thick},
  net/.style={rectangle, draw, rounded corners=4pt, fill=blue!20, minimum width=2.2cm, minimum height=1.2cm, align=center, font=\small\sffamily, draw=blue!80!black, thick},
  aux/.style={rectangle, draw, rounded corners=4pt, fill=gray!20, minimum width=2.2cm, minimum height=1.2cm, align=center, font=\small\sffamily, draw=gray!80!black, thick},
  loss/.style={diamond, draw, fill=red!20, minimum width=1.6cm, minimum height=1.6cm, align=center, font=\scriptsize\sffamily, draw=red!80!black, thick},
  out/.style={rectangle, draw, fill=green!20, minimum width=1.8cm, minimum height=0.9cm, align=center, font=\small\sffamily\bfseries, rounded corners=2pt, draw=green!80!black, thick}
]

% ================= 节点定义 =================
% 输入数据
\node[data] (input_sais) at (0, 5) {光场 SAIs\\{\scriptsize $(N^2, 3, H, W)$}};
\node[data] (input_ic) at (0, 0) {中心视角 $I_c$\\{\scriptsize $(3, H, W)$}};

% 特征提取分支
\node[core] (ang_fft) at (2.2, 7.5) {角度FFT特征\\提取模块};
\node[aux] (shallow_cnn) at (2.2, 3.0) {浅层CNN};

\node[data] (freq_feat) at (4.4, 7.5) {频域特征 $F$\\{\scriptsize $(H,W,C,F_u,F_v)$}};
\node[data] (spatial_feat) at (4.4, 3.0) {空间特征\\{\scriptsize $(H,W,C_{spatial})$}};

% 双掩码生成
\node[core] (dual_mask) at (6.6, 5.25) {双掩码\\生成模块};
\node[loss] (mutual_loss) at (6.6, 1.5) {互斥损失};

\node[data] (mask_md) at (8.8, 6.5) {漫反射掩码 $M_d$\\{\scriptsize $(H,W,1)$}};
\node[data] (mask_ms) at (8.8, 4.0) {高光掩码 $M_s$\\{\scriptsize $(H,W,1)$}};

% 物理分离与深度估计
\node[net] (phys_sep) at (11.0, 5.25) {物理反射\\分离模块};
\node[core] (unified_depth) at (13.2, 5.25) {统一深度\\估计模块};
\node[data] (d_init) at (15.4, 5.25) {初始深度 $D_{init}$\\{\scriptsize $(H,W,1)$}};

% 深度细化与输出
\node[net] (depth_refine) at (15.4, 1.5) {深度细化模块};
\node[out] (d_final) at (17.6, 1.5) {最终深度 $D_{final}$\\{\scriptsize $(H,W,1)$}};

% ================= 数据流连接 =================
% 1. 输入到特征提取
\draw[arrow] (input_sais.east) -- ++(0.9,0) |- (ang_fft.west);
\draw[arrow] (input_sais.east) -- ++(0.9,0) |- (shallow_cnn.west);

% 2. 特征提取到特征图
\draw[arrow] (ang_fft) -- node[above, font=\scriptsize] {2D FFT} (freq_feat);
\draw[arrow] (shallow_cnn) -- node[above, font=\scriptsize] {Conv} (spatial_feat);

% 3. 特征图到双掩码生成
\draw[arrow] (freq_feat.east) -- ++(0.9,0) coordinate (c1) |- (dual_mask.west);
\node[above, font=\scriptsize] at (c1) {Pool};
\draw[arrow] (spatial_feat.east) -- ++(0.9,0) coordinate (c2) |- (dual_mask.west);
\node[below, font=\scriptsize] at (c2) {Concat};

% 4. 双掩码生成到掩码及损失
\draw[arrow_dash] (dual_mask.south) -- node[right, font=\scriptsize] {约束} (mutual_loss.north);
\draw[arrow] (dual_mask.east) -- ++(0.9,0) coordinate (c3) |- (mask_md.west);
\draw[arrow] (dual_mask.east) -- ++(0.9,0) coordinate (c4) |- (mask_ms.west);

% 5. 原始SAIs到物理分离 (底部绕线)
\draw[arrow] (input_sais.south) -- ++(0,-7.5) coordinate (c_sais) -| (phys_sep.south);
\node[right, font=\scriptsize] at (c_sais) {原始 SAIs};

% 6. 掩码到物理分离
\draw[arrow] (mask_md.east) -- ++(0.9,0) coordinate (c5) |- (phys_sep.west);
\node[above, font=\scriptsize] at (c5) {$M_d$};
\draw[arrow] (mask_ms.east) -- ++(0.9,0) coordinate (c6) |- (phys_sep.west);
\node[below, font=\scriptsize] at (c6) {$M_s$};

% 7. 物理分离到统一深度估计
\draw[arrow] (phys_sep) -- node[above, font=\scriptsize] {$L_d, L_s$} (unified_depth);

% 8. 掩码到统一深度估计 (辅助信息流)
\draw[arrow_dash] (mask_md.east) -- ++(2.8,0) coordinate (c7) |- (unified_depth.north);
\draw[arrow_dash] (mask_ms.east) -- ++(2.8,0) coordinate (c8) |- (unified_depth.south);

% 9. 统一深度估计到初始深度
\draw[arrow] (unified_depth) -- node[above, font=\scriptsize] {Soft-Argmin} (d_init);

% 10. 初始深度与中心视角到深度细化
\draw[arrow] (d_init.south) -- ++(0,-1.5) -| (depth_refine.north);
\draw[arrow] (input_ic.east) -- ++(13.0,0) |- (depth_refine.west);

% 11. 深度细化到最终输出
\draw[arrow] (depth_refine) -- node[above, font=\scriptsize] {引导滤波} (d_final);

% ================= 图例 =================
\begin{scope}[shift={(0, -3.5)}]
  \node[fill=orange!30, draw, rounded corners=2pt, minimum width=0.5cm, minimum height=0.3cm, draw=orange!80!black, thick] (leg1) at (0, 0) {};
  \node[right, font=\scriptsize\sffamily] at (leg1.east) {核心创新模块};

  \node[fill=blue!20, draw, rounded corners=2pt, minimum width=0.5cm, minimum height=0.3cm, draw=blue!80!black, thick] (leg2) at (3.5, 0) {};
  \node[right, font=\scriptsize\sffamily] at (leg2.east) {网络处理模块};

  \node[fill=gray!20, draw, rounded corners=2pt, minimum width=0.5cm, minimum height=0.3cm, draw=gray!80!black, thick] (leg3) at (7.0, 0) {};
  \node[right, font=\scriptsize\sffamily] at (leg3.east) {辅助模块};

  \node[fill=red!20, draw, diamond, minimum width=0.4cm, minimum height=0.4cm, draw=red!80!black, thick] (leg4) at (10.5, 0) {};
  \node[right, font=\scriptsize\sffamily] at (leg4.east) {损失/约束};
  
  \node[fill=blue!10, draw, rounded corners=2pt, minimum width=0.5cm, minimum height=0.3cm, draw=blue!60!black] (leg5) at (14.0, 0) {};
  \node[right, font=\scriptsize\sffamily] at (leg5.east) {张量/特征图};
\end{scope}

\end{tikzpicture}